% paper/sections/12a_force_balance.tex
%
% §12.1  力学的バランス検証（Force Balance Verification）
% ═══════════════════════════════════════════════════════════════════════════════
\subsection{力学的バランス検証}
\label{sec:force_balance}
% ═══════════════════════════════════════════════════════════════════════════════

NS ソルバの最も基本的な物理的整合性は，
\textbf{静止流体における力の釣合が数値的に維持されること}である．
本節では静水圧均衡と表面張力均衡の 2 種類のテストにより検証する．

% ─────────────────────────────────────────────────────────────────────────────
\subsubsection{静水圧均衡}
\label{sec:hydrostatic}
% ─────────────────────────────────────────────────────────────────────────────

\noindent\textbf{目的：}
Predictor--PPE--Corrector パイプラインが重力と圧力勾配の離散的釣合を
厳密に維持することを確認する．

\noindent\textbf{評価対象：}
NS ソルバ全体の離散力平衡（CCD 微分 + PPE 投影 + 速度補正）．

\noindent\textbf{手順とパラメタ：}
\begin{itemize}[noitemsep]
  \item 計算領域 $[0,1]^2$，単相流，壁面境界条件，重力 $\bm{g} = (0, -1)^\top$
  \item 初期条件 $\bu = \bm{0}$，$p = \rho g (1-y)$（静水圧から出発）
  \item 格子 $N = 32, 64, 128$，100 ステップ実行
  \item 合格基準 $\|\bu\|_\infty < 10^{-6}$（$N \geq 64$ では $10^{-9}$ 以下を期待）
\end{itemize}

\noindent\textbf{結果：}

\begin{table}[htbp]
\centering
\caption{静水圧均衡テスト：100 ステップ後の速度・圧力誤差}
\label{tab:hydrostatic}
\renewcommand{\arraystretch}{1.3}
\begin{tabular}{rrrc}
\toprule
$N$ & $\|\bu\|_\infty$ & $|p - p_{\mathrm{hydro}}|_\infty$ & 判定 \\
\midrule
32  & $2.45 \times 10^{-6}$ & $7.08 \times 10^{-5}$ & 合格 \\
64  & $2.31 \times 10^{-10}$ & $7.02 \times 10^{-9}$ & 合格 \\
128 & $6.78 \times 10^{-10}$ & $2.20 \times 10^{-8}$ & 合格 \\
\bottomrule
\end{tabular}
\end{table}

\noindent\textbf{結論：}
$N = 64$ で速度は $\Ord{10^{-10}}$ に収まり，
IPC による重力と圧力勾配の離散的釣合が 100 ステップにわたり維持された．
$N = 32$ では CCD 壁面境界スキームの低次精度（$\Ord{h^5}$）に起因する初期誤差が IPC を通じて蓄積するため
$\Ord{10^{-6}}$ に留まるが，格子精緻化により急速に減衰する．

% ─────────────────────────────────────────────────────────────────────────────
\subsubsection{表面張力均衡と寄生流れ}
\label{sec:bf_static_droplet}
% ─────────────────────────────────────────────────────────────────────────────

\noindent\textbf{目的：}
静止液滴における Laplace 圧均衡（$\Delta p = \sigma/R$）を検証し，
CSF + 変密度 PPE の離散化が正しい圧力跳びを再現することを確認する．

\noindent\textbf{評価対象：}
CCD 曲率 + CSF 体積力 + 変密度 PPE（$\bnabla\cdot({\rho^{-1}}\bnabla p)$）の
単ステップ投影による Laplace 圧精度．

\noindent\textbf{手順とパラメタ：}
\begin{itemize}[noitemsep]
  \item 計算領域 $[0,1]^2$，壁面境界条件，重力なし
  \item 静止液滴 $R = 0.25$，$\rho_l/\rho_g = 1000$，$\mathit{We} = 1$
  \item CSF 体積力 $\bm{f}_\sigma = (1/\mathit{We})\,\kappa^*\,\bnabla\psi$（無次元形式では $\sigma$ は $We$ に吸収済み）
  \item 曲率 $\kappa^*$ は InterfaceLimitedFilter（$C = 0.05$）適用後の値：
        $\kappa^* = \kappa + C\,h^2\,w(\psi)\,\nabla^2\kappa$，$w(\psi) = 4\psi(1-\psi)$
  \item 単ステップ投影：$\bu^* = \Delta t\,\bm{f}_\sigma / \rho$ → 変密度 PPE → Corrector
  \item PPE 右辺の発散は CCD 微分（$\Ord{h^6}$）で計算
  \item 格子 $N = 32, 64, 128, 256$，厳密解 $\Delta p = \sigma / R = 4.0$
\end{itemize}

\noindent\textbf{結果：}

\begin{table}[htbp]
\centering
\caption{静止液滴 Laplace 圧テスト：変密度 PPE + CSF（$\rho_l/\rho_g = 1000$）}
\label{tab:bf_laplace}
\renewcommand{\arraystretch}{1.3}
\begin{tabular}{rrrrr}
\toprule
$N$ & $\Delta p$ & 相対誤差 & $\|\bu_{\mathrm{para}}\|_\infty$ & $\Delta p$ 次数 \\
\midrule
 32 & $7.72$  & $9.3 \times 10^{-1}$ & $1.7 \times 10^{-3}$ & --- \\
 64 & $4.08$  & $2.1 \times 10^{-2}$ & $7.6 \times 10^{-5}$ & $5.47$ \\
128 & $3.98$  & $6.1 \times 10^{-3}$ & $3.0 \times 10^{-5}$ & $1.79$ \\
256 & $3.99$  & $2.2 \times 10^{-3}$ & $2.3 \times 10^{-4}$ & $1.45$ \\
\bottomrule
\end{tabular}
\end{table}

\noindent\textbf{結論：}
$N = 64, 128, 256$ の 3 水準で $\Delta p$ の相対誤差は単調減少し，
収束次数は $5.47 \to 1.79 \to 1.45$ と漸近的に CSF モデル誤差 $\Ord{h^2}$ に近づく%
\footnote{$N = 32$ では $\Delta p = 7.72$（負の値から正に反転後もなお過大）であり，
$R/h = 8$ は界面幅 $\varepsilon$ が液滴半径 $R$ に対して過大で
$\delta_\varepsilon$ 正則化の仮定が破綻するため漸近領域外にある．}．
$N = 256$ で $\Delta p = 3.99$（相対誤差 $0.22\%$）であり，
CCD-PPE の高精度離散化が CSF $\Ord{h^2}$ の枠内で十分に機能している．

図~\ref{fig:ch12_droplet_fields} に $N = 64$ における 2 次元場を示す．
上段は圧力場（Laplace ジャンプ），寄生電流の空間分布，
および寄生流れ振幅の時間履歴である．
下段には HFE 曲率フィルタ適用後の渦度場，流線（速度大きさで着色），
および速度ベクトルを示す．
渦度は界面上で 4 極パターンを呈し，
流線は 4 つの循環セルで構成される寄生電流の構造を可視化している．

\begin{figure}[p]
\centering
\includegraphics[width=\linewidth,height=0.85\textheight,keepaspectratio]{figures/ch12_droplet_fields}
\caption{静止液滴 2 次元場（$N = 64$，$\rho_l/\rho_g = 2$，$\mathit{We} = 10$，200 ステップ）．
上段左：圧力場 $p(x,y)$，上段中：寄生電流 $\|\bu\|(x,y)$，上段右：$\|\bu\|_\infty$ 時間履歴．
下段左：渦度 $\omega(x,y)$，下段中：流線（速度で着色），下段右：速度ベクトル．
黒破線および赤線は界面 $\phi = 0$．}
\label{fig:ch12_droplet_fields}
\end{figure}

寄生流れ振幅は $N = 64 \to 128$ で減少するが，$N = 256$ で増大する．
これは FD PPE（$\Ord{h^2}$）と CCD 圧力勾配（$\Ord{h^6}$）の間の
投影不整合が高解像度で顕在化するためである
（§~\ref{sec:ppe_discretization_choice}，付録参照）．
この不整合が多ステップ時間積分において累積するかを
§~\ref{sec:advection_pressure_coupling} の二相流時間積分テストで確認する．

\paragraph{CCD 勾配 vs FD 勾配の寄生流れ比較}

§~\ref{sec:ccd_balanced_force} の Balanced-Force 条件——表面張力と圧力勾配に
\textbf{同一の}離散勾配演算子を用いる——の効果を定量化するため，
静的液滴テスト（$R = 0.25$，$\rho_l/\rho_g = 1$，$\mathit{We} = 1$，200 ステップ）において
CCD 勾配と 2 次 FD 勾配を切り替えて寄生流れを比較した．

$N = 64$ での結果：CCD 勾配では $\|\bu_{\mathrm{para}}\|_\infty = 4.5 \times 10^{-3}$，
FD 勾配では $\|\bu_{\mathrm{para}}\|_\infty = 5.0 \times 10^{-2}$ であり，
FD は CCD の\textbf{約 11 倍}の寄生流れを生じた．
CCD 勾配では $N = 32 \to 64$ で寄生流れが $\Ord{h^{5.3}}$ で減少した一方，
FD 勾配の収束次数は $\Ord{h^{2.2}}$ にとどまった．
この差は，CCD が $\bnabla p$ と $\kappa\bnabla H$ の離散化誤差を $\Ord{h^6}$ で
相殺するのに対し，FD では $\Ord{h^2}$ の不整合が残ることに起因する
（§~\ref{sec:ccd_balanced_force}）．

% C/RC 補正の効果 → §12.4.1 の静止液滴テーブルに統合
